%%%%%%%%%%%%%%%%
%%% gen 25 %%%
%%%%%%%%%%%%%%%%

\Chapter{25}

\Verse{1}
{
	\hE{וַיֹּסֶף אַבְרָהָם וַיִּקַּח אִשָּׁה וּשְׁמָהּ קְטוּרָה׃}
}
{
	\eE{
		And Abraham went on and took a wife, and her name was
		Keturah,
	}
}
{
%	\footnote{
%		Keturah. The most ignored significant person in the Torah. Rashi follows
%		an old rabbinic idea that she is Hagar. But there is no basis for this
%		in the text, and other traditional commentators reject it (Ibn Ezra,
%		Ramban, Rashbam). Keturah is the mother of tribes located along the
%		route of incense trade, and her name is related to the word for incense.
%		Notably, the Midianites are among the children of Abraham and Keturah,
%		and the influence of the Midianite priest Jethro on Moses, his
%		son-in-law, is understood to be substantial. And the line of Levites who
%		are descended from Moses thus—alone among the Israelites—derive from
%		Abraham through both Sarah and Keturah.
%	}
}

\Verse{2}
{
	\hE{וַתֵּלֶד לוֹ אֶת זִמְרָן וְאֶת יקְשָׁן וְאֶת מְדָן וְאֶת מִדְיָן וְאֶת יִשְׁבָּק וְאֶת שׁוּחַ׃}
}
{
	\eE{
		and she gave birth to Zimran and Jokshan and Medan and
		Midian and Ishbak and Shuah for him.
	}
}
{
}

\Verse{3}
{
	\hE{וְיקְשָׁן יָלַד אֶת שְׁבָא וְאֶת דְּדָן וּבְנֵי דְדָן הָיוּ אַשּׁוּרִם וּלְטוּשִׁם וּלְאֻמִּים׃}
}
{
	\eE{
		And Jokshan fathered Sheba and Dedan. And the children of
		Dedan were Ashurim and Letushim and Leummim.
	}
}
{
}

\Verse{4}
{
	\hE{וּבְנֵי מִדְיָן עֵיפָה וָעֵפֶר וַחֲנֹךְ וַאֲבִידָע וְאֶלְדָּעָה כּל אֵלֶּה בְּנֵי קְטוּרָה׃}
}
{
	\eE{
		And the children of Midian: Ephah and Epher and Hanoch and
		Abida and Eldaah. All these were children of Keturah.
	}
}
{
}

\Verse{5}
{
	\hRJE{וַיִּתֵּן אַבְרָהָם אֶת כּל אֲשֶׁר לוֹ לְיִצְחָק׃}
}
{
	\eRJE{And Abraham gave everything that he had to Isaac.}
}
{
}

\Verse{6}
{
	\hRJE{וְלִבְנֵי הַפִּילַגְשִׁים אֲשֶׁר לְאַבְרָהָם נָתַן אַבְרָהָם מַתָּנֹת וַיְשַׁלְּחֵם מֵעַל יִצְחָק בְּנוֹ בְּעוֹדֶנּוּ חַי קֵדְמָה אֶל אֶרֶץ קֶדֶם׃}
}
{
	\eRJE{
		And Abraham gave gifts to the children of the concubines
		that Abraham had, and he sent them away from Isaac, his
		son, while he was still living: east, to the land of the
		East.
	}
}
{
}

\Verse{7}
{
	\hP{וְאֵלֶּה יְמֵי שְׁנֵי חַיֵּי אַבְרָהָם אֲשֶׁר חָי מְאַת שָׁנָה וְשִׁבְעִים שָׁנָה וְחָמֵשׁ שָׁנִים׃}
}
{
	\eP{
		And these are the days of the years of Abraham’s life that
		he lived: a hundred years and seventy years and five
		years.
	}
}
{
}

\Verse{8}
{
	\hP{וַיִּגְוַע וַיָּמת אַבְרָהָם בְּשֵׂיבָה טוֹבָה זָקֵן וְשָׂבֵעַ וַיֵּאָסֶף אֶל עַמָּיו׃}
}
{
	\eP{
		And he expired. And Abraham died at a good old age, old
		and full, and was gathered to his people.
	}
}
{
}

\Verse{9}
{
	\hP{וַיִּקְבְּרוּ אֹתוֹ יִצְחָק וְיִשְׁמָעֵאל בָּנָיו אֶל מְעָרַת הַמַּכְפֵּלָה אֶל שְׂדֵה עֶפְרֹן בֶּן צֹחַר הַחִתִּי אֲשֶׁר עַל פְּנֵי מַמְרֵא׃}
}
{
	\eP{
		And Isaac and Ishmael, his sons, buried him at the cave of
		Machpelah at the field of Ephron, son of Zohar, the
		Hittite, facing Mamre,
	}
}
{
}

\Verse{10}
{
	\hP{הַשָּׂדֶה אֲשֶׁר קָנָה אַבְרָהָם מֵאֵת בְּנֵי חֵת שָׁמָּה קֻבַּר אַבְרָהָם וְשָׂרָה אִשְׁתּוֹ׃}
}
{
	\eP{
		the field that Abraham bought from the children of Heth.
		Abraham was buried there—and Sarah, his wife.
	}
}
{
}

\Verse{11}
{
	\hP{וַיְהִי אַחֲרֵי מוֹת אַבְרָהָם וַיְבָרֶךְ אֱלֹהִים אֶת יִצְחָק בְּנוֹ וַיֵּשֶׁב יִצְחָק עִם בְּאֵר לַחַי רֹאִי׃}
}
{
	\eP{
		And it was after Abraham’s death, and {\God} blessed Isaac,
		his son. And Isaac lived at the well Lahai-roi.
	}
}
{
%	\footnote{
%		Isaac lived at the well Lahai-roi. He lives at the place where Hagar was
%		told by an angel that she would give birth to Ishmael, Isaac’s brother
%		(Gen 16:14). Ishmael remains significant in the background. It is as if
%		Isaac recognizes that his miraculous birth and his being spared from
%		sacrifice have given him everything that otherwise would have been
%		Ishmael’s. Living at Lahai-roi, he is reminded that Ishmael’s life, too,
%		was divinely determined.
%	}
}

\Verse{12}
{
	\hR{וְאֵלֶּה תֹּלְדֹת יִשְׁמָעֵאל בֶּן אַבְרָהָם אֲשֶׁר יָלְדָה הָגָר הַמִּצְרִית שִׁפְחַת שָׂרָה לְאַבְרָהָם׃}
}
{
	\eR{
		And these are the records of Ishmael, son of Abraham, to
		whom Hagar, the Egyptian, Sarah’s maid, gave birth for
		Abraham:
	}
}
{
}

\Verse{13}
{
	\hP{וְאֵלֶּה שְׁמוֹת בְּנֵי יִשְׁמָעֵאל בִּשְׁמֹתָם לְתוֹלְדֹתָם בְּכֹר יִשְׁמָעֵאל נְבָיֹת וְקֵדָר וְאַדְבְּאֵל וּמִבְשָׂם׃}
}
{
	\eP{
		And these are the names of Ishmael’s sons, by their names,
		according to their records: Ishmael’s firstborn was
		Nebaioth, and Kedar and Adbeel and Mibsam
	}
}
{
}

\Verse{14}
{
	\hP{וּמִשְׁמָע וְדוּמָה וּמַשָּׂא׃}
}
{
	\eP{and Mishma and Dumah and Massa,}
}
{
}

\Verse{15}
{
	\hP{חֲדַד וְתֵימָא יְטוּר נָפִישׁ וָקֵדְמָה׃}
}
{
	\eP{Hadar and Tema, Jetur, Naphish and Kedmah.}
}
{
}

\Verse{16}
{
	\hP{אֵלֶּה הֵם בְּנֵי יִשְׁמָעֵאל וְאֵלֶּה שְׁמֹתָם בְּחַצְרֵיהֶם וּבְטִירֹתָם שְׁנֵים עָשָׂר נְשִׂיאִם לְאֻמֹּתָם׃}
}
{
	\eP{
		These are they, Ishmael’s sons, and these are their names
		in their settlements and in their encampments, twelve
		chieftains by their clans.
	}
}
{
}

\Verse{17}
{
	\hP{וְאֵלֶּה שְׁנֵי חַיֵּי יִשְׁמָעֵאל מְאַת שָׁנָה וּשְׁלֹשִׁים שָׁנָה וְשֶׁבַע שָׁנִים וַיִּגְוַע וַיָּמת וַיֵּאָסֶף אֶל עַמָּיו׃}
}
{
	\eP{
		And these are the years of Ishmael’s life: a hundred years
		and thirty years and seven years. And he expired and died
		and was gathered to his people.
	}
}
{
}

\Verse{18}
{
	\hP{וַיִּשְׁכְּנוּ מֵחֲוִילָה עַד שׁוּר אֲשֶׁר עַל פְּנֵי מִצְרַיִם בֹּאֲכָה אַשּׁוּרָה עַל פְּנֵי כל אֶחָיו נָפָל׃}
}
{
	\eP{
		And they tented from Havilah to Shur, which faces Egypt,
		as you come toward Asshur. He fell facing all of his
		brothers.
	}
}
{
%	\footnote{
%		He fell facing all of his brothers. Meaning: his family’s places of
%		residence were located among the territories of all their kindred
%		peoples (presumably Israelites, Moabites, Ammonites, Edomites).
%	}
}

\Verse{19}
{
	\hR{וְאֵלֶּה תּוֹלְדֹת יִצְחָק בֶּן אַבְרָהָם אַבְרָהָם הוֹלִיד אֶת יִצְחָק׃}
}
{
	\eR{
		RECORDS And these are the records of Isaac, son of
		Abraham: Abraham had fathered Isaac.
	}
}
{
%	\footnote{
%		these are the records of Isaac. Isaac’s life is different from his
%		father Abraham’s and his son Jacob’s. Unlike them: He never leaves the
%		promised land. He never fights. He has only one wife. His name is not
%		changed. He is pictured going to meditate. He lives the longest. There
%		are fewer stories about him. He is a relatively passive figure, acted on
%		by his father and his wife and his son. Why is he different? Presumably
%		because he has once lain on an altar as a sacrifice to God. Even though
%		the sacrifice is not consummated, his life is now consecrated, and his
%		life is distinct after this. The actions that parents take regarding
%		their children impact on their children’s lives ever after.
%	}
}

\Verse{20}
{
	\hP{וַיְהִי יִצְחָק בֶּן אַרְבָּעִים שָׁנָה בְּקַחְתּוֹ אֶת רִבְקָה בַּת בְּתוּאֵל הָאֲרַמִּי מִפַּדַּן אֲרָם אֲחוֹת לָבָן הָאֲרַמִּי לוֹ לְאִשָּׁה׃}
}
{
	\eP{
		And Isaac was forty years old when he took Rebekah,
		daughter of Bethuel, the Aramean, from Paddan Aram, sister
		of Laban, the Aramean, to him as a wife.
	}
}
{
}

\Verse{21}
{
	\hJ{וַיֶּעְתַּר יִצְחָק לַיהֹוָה לְנֹכַח אִשְׁתּוֹ כִּי עֲקָרָה הִוא וַיֵּעָתֶר לוֹ יְהֹוָה וַתַּהַר רִבְקָה אִשְׁתּוֹ׃}
}
{
	\eJ{
		And Isaac prayed to YHWH for his wife because she was
		infertile, and YHWH was prevailed upon by him, and his
		wife Rebekah became pregnant.
	}
}
{
}

\Verse{22}
{
	\hJ{וַיִּתְרֹצְצוּ הַבָּנִים בְּקִרְבָּהּ וַתֹּאמֶר אִם כֵּן לָמָּה זֶּה אָנֹכִי וַתֵּלֶךְ לִדְרֹשׁ אֶת יְהֹוָה׃}
}
{
	\eJ{
		And the children struggled inside her, and she said, “If
		it’s like this, why do I exist?” And she went to inquire
		of YHWH.
	}
}
{
}

\Verse{23}
{
	\hJ{וַיֹּאמֶר יְהֹוָה לָהּ שְׁנֵי (גיים) [גוֹיִם] בְּבִטְנֵךְ וּשְׁנֵי לְאֻמִּים מִמֵּעַיִךְ יִפָּרֵדוּ וּלְאֹם מִלְאֹם יֶאֱמָץ וְרַב יַעֲבֹד צָעִיר׃}
}
{
	\eJ{
		And YHWH said to her, Two nations are in your womb, and
		two peoples will be dispersed from your insides, and one
		people will be mightier than the other people, and the
		older the younger will serve.
	}
}
{
%	\footnote{
%		YHWH said to her. Rebekah hears the voice of God before Isaac does. (God
%		first speaks to Isaac in the next chapter, 26:2.) In matters of
%		revelation, man is not more important than woman. The message is what is
%		important. Rebekah hears first because her need arises first. She is
%		troubled because of the struggling twins in her womb. Footnote 25:23.
%		the older the younger will serve. As I explained in The Hidden Face of
%		God, people have usually taken this to mean that YHWH tells Rebekah that
%		her younger son, Jacob, will dominate her older one, Esau. Some have
%		thought, therefore, that Rebekah is not really manipulating the
%		succession when she sends Jacob to pose as Esau. Rather, she is simply
%		fulfilling God’s will. The decision of who is number-one-son thus is
%		still God’s. However, this understanding is based on a misunderstanding
%		of the subtle, exquisitely ambiguous biblical wording. The text does not
%		in fact say that the elder son will serve the younger son. In biblical
%		Hebrew, the subject may either precede or follow the verb, and the
%		object likewise may either precede or follow the verb. What that means
%		is that sometimes it is impossible to tell which word in a biblical
%		verse is the subject and which is the object, especially if the verse is
%		in poetry. That is the case in this oracle to Rebekah, which is in
%		poetry. It can mean: “the elder will serve the younger” But it equally
%		can mean: “the elder, the younger will serve” Like the Delphic oracles
%		in Greece, this prediction contains two opposite meanings, and thus the
%		person who receives it—Rebekah—can hear whatever she wants (consciously
%		or subconsciously) to hear. It can be understood to mean that Jacob will
%		serve Esau or that Esau will serve Jacob. I learned this from my senior
%		colleague, David Noel Freedman.
%	}
}

\Verse{24}
{
	\hJ{וַיִּמְלְאוּ יָמֶיהָ לָלֶדֶת וְהִנֵּה תוֹמִם בְּבִטְנָהּ׃}
}
{
	\eJ{
		And her days to give birth were completed, and here were
		twins in her womb.
	}
}
{
}

\Verse{25}
{
	\hJ{וַיֵּצֵא הָרִאשׁוֹן אַדְמוֹנִי כֻּלּוֹ כְּאַדֶּרֶת שֵׂעָר וַיִּקְרְאוּ שְׁמוֹ עֵשָׂו׃}
}
{
	\eJ{
		And the first came out all ruddy, like a hairy robe, and
		they called his name Esau.
	}
}
{
}

\Verse{26}
{
	\hJ{וְאַחֲרֵי כֵן יָצָא אָחִיו וְיָדוֹ אֹחֶזֶת בַּעֲקֵב עֵשָׂו וַיִּקְרָא שְׁמוֹ יַעֲקֹב וְיִצְחָק בֶּן שִׁשִּׁים שָׁנָה בְּלֶדֶת אֹתָם׃}
}
{
	\eJ{
		And after that his brother came out, and his hand was
		holding Esau’s heel, and he called his name Jacob. And
		Isaac was sixty years old at their birth.
	}
}
{
}

\Verse{27}
{
	\hJ{וַיִּגְדְּלוּ הַנְּעָרִים וַיְהִי עֵשָׂו אִישׁ יֹדֵעַ צַיִד אִישׁ שָׂדֶה וְיַעֲקֹב אִישׁ תָּם יֹשֵׁב אֹהָלִים׃}
}
{
	\eJ{
		And the boys grew up, and Esau was a man who knew hunting,
		a man of the field, and Jacob was a simple man, living in
		tents.
	}
}
{
%	\footnote{
%		Esau. Numerous attempts have been made to denigrate Esau in midrash and
%		even in current biblical interpretation. It is not justified according
%		to the text. The motive is understandable: Jacob’s behavior in the
%		matters of the birthright and the blessing is an embarrassment. Even
%		small children express surprise at what this patriarch does to his
%		brother—and his father. The Torah itself neither denigrates Esau nor
%		justifies Jacob. On the contrary, one of the great qualities of the
%		Tanak is precisely that none of its heroes is perfect. Moreover, Jacob
%		suffers some corresponding recompense for every deception he commits, as
%		do Laban, Rachel, Reuben, Simeon, Levi, Judah, and Joseph’s other
%		brothers. (See the relevant comments below.) The lesson, therefore, is
%		not only that people are not perfect, but also that acts of deception
%		can fester in a family, and that they have consequences. Those who try
%		to excuse Jacob and turn Esau into a villain are unfortunately hiding
%		these essential lessons of the Torah.
%	}
}

\Verse{28}
{
	\hJ{וַיֶּאֱהַב יִצְחָק אֶת עֵשָׂו כִּי צַיִד בְּפִיו וְרִבְקָה אֹהֶבֶת אֶת יַעֲקֹב׃}
}
{
	\eJ{
		And Isaac loved Esau because he put game meat in his
		mouth, and Rebekah loved Jacob.
	}
}
{
}

\Verse{29}
{
	\hJ{וַיָּזֶד יַעֲקֹב נָזִיד וַיָּבֹא עֵשָׂו מִן הַשָּׂדֶה וְהוּא עָיֵף׃}
}
{
	\eJ{
		And Jacob made a stew, and Esau came from the field, and
		he was exhausted.
	}
}
{
}

\Verse{30}
{
	\hJ{וַיֹּאמֶר עֵשָׂו אֶל יַעֲקֹב הַלְעִיטֵנִי נָא מִן הָאָדֹם הָאָדֹם הַזֶּה כִּי עָיֵף אָנֹכִי עַל כֵּן קָרָא שְׁמוֹ אֱדוֹם׃}
}
{
	\eJ{
		And Esau said to Jacob, “Will you feed me some of the red
		stuff, this red stuff, because I’m exhausted.” (On account
		of this he called his name “Edom.”)
	}
}
{
}

\Verse{31}
{
	\hJ{וַיֹּאמֶר יַעֲקֹב מִכְרָה כַיּוֹם אֶת בְּכֹרָתְךָ לִי׃}
}
{
	\eJ{And Jacob said, “Sell your birthright to me, today.”}
}
{
}

\Verse{32}
{
	\hJ{וַיֹּאמֶר עֵשָׂו הִנֵּה אָנֹכִי הוֹלֵךְ לָמוּת וְלָמָּה זֶּה לִי בְּכֹרָה׃}
}
{
	\eJ{
		And Esau said, “Here I’m going to die, and what use is
		this, that I have a birthright?”
	}
}
{
}

\Verse{33}
{
	\hJ{וַיֹּאמֶר יַעֲקֹב הִשָּׁבְעָה לִּי כַּיּוֹם וַיִּשָּׁבַע לוֹ וַיִּמְכֹּר אֶת בְּכֹרָתוֹ לְיַעֲקֹב׃}
}
{
	\eJ{
		And Jacob said, “Swear to me, today.” And he swore to him
		and sold his birthright to Jacob. And Jacob gave bread and
		lentil stew to Esau. And he ate and drank and got up and
		went. And Esau disdained the birthright.
	}
}
{
}

\Verse{34}
{
	\hOther{וְיַעֲקֹב נָתַן לְעֵשָׂו לֶחֶם וּנְזִיד עֲדָשִׁים וַיֹּאכַל וַיֵּשְׁתְּ וַיָּקם וַיֵּלַךְ וַיִּבֶז עֵשָׂו אֶת הַבְּכֹרָה׃}
}
{
	\eOther{}
}
{
}
